\TUKEsection{Pamäte}{Flash, SRAM a EEPROM}

\begin{frame}{Základné pojmy}
\begin{columns}[T,onlytextwidth]
\column{0.48\textwidth}
\begin{tukebox}{Hardvérový pohľad}
\begin{itemize}
    \item periférny blok,
    \item konfiguračný register,
    \item stavový príznak,
    \item prerušenie alebo polling.
\end{itemize}
\end{tukebox}
\column{0.48\textwidth}
\begin{tukebox}{Programátorský pohľad}
\begin{itemize}
    \item nastavenie bitu,
    \item čítanie registra,
    \item maskovanie,
    \item testovanie stavu.
\end{itemize}
\end{tukebox}
\end{columns}
\end{frame}

\begin{frame}{Typický postup práce s perifériou}
\begin{enumerate}
    \item nájsť príslušnú kapitolu v datasheete,
    \item identifikovať registre a bity,
    \item vypočítať potrebné parametre,
    \item nakonfigurovať perifériu,
    \item overiť funkciu meraním alebo diagnostickým výpisom.
\end{enumerate}
\end{frame}
